% !TeX root = ../../progress-sem1.tex
% Section 3.5: UI prototype, split like the example
% (การออกแบบ UI / prototype ของผู้ใช้แต่ละกลุ่ม).
% Input from chapters/03-methodology.tex.

\section{การออกแบบ User Interface Prototype}\label{sec:ui-prototype}
\thaijustify{
  ทีมออกแบบหน้าจอแบบ high-fidelity ด้วย Figma ก่อนพัฒนา โดยแบ่งตามกลุ่มผู้ใช้ของแพลตฟอร์ม
}

\subsection{การออกแบบ UI (User Interface)}
\thaijustify{
  หน้าจอของแพลตฟอร์มใช้ design system ที่สร้างบน shadcn/ui และ Tailwind CSS ไอคอนใช้ชุด Hugeicons ทั้งระบบ
  \todo{โทนสี ตัวอักษร ขนาด spacing และ component หลักใน Figma พร้อมภาพ style guide}
}

\subsection{User Interface Prototype ของ Designer}
\thaijustify{
  \todo{ภาพ prototype ของ Designer: dashboard, การสร้างโปรเจกต์และ design intent, user story, การตรวจ prompt, generation trace และหน้า refine}
}

\subsection{User Interface Prototype ของ Usability Tester}
\thaijustify{
  \todo{ภาพ prototype ของผู้เข้าร่วมทดสอบ: หน้าเข้าร่วมเซสชัน, การให้ความยินยอม (PDPA), การทำภารกิจ และแบบฟอร์มข้อเสนอแนะ}
}

\subsection{User Interface Prototype ของ Client, Developer และ Project Manager}
\thaijustify{
  \todo{ภาพ prototype ของมุมมองผู้มีส่วนได้ส่วนเสีย: สถานะโปรเจกต์, handoff และรายงาน}
}
